\Conclusion % заключение к отчёту

В результате выполнения выпускной квалификационной работы спроектирован и реализован программный комплекс для централизованного управления и мониторинга распределённой инфраструктуры имитационных сервисов. Все поставленные задачи решены в полном объёме.

В ходе анализа предметной области установлено, что существующие зарубежные решения (WireMock, MockServer, Mountebank) не обеспечивают совместимость с отечественными сертифицированными операционными системами семейства Astra Linux и не поддерживают централизованное управление распределённой инфраструктурой Java-процессов без установки агентного программного обеспечения на целевые серверы. Сравнительный анализ инструментальных средств позволил обосновать технологический стек: Python / FastAPI / Redis, -- оптимально сочетающий асинхронную модель обработки запросов с минимальной задержкой доступа к данным конфигурации.

При проектировании архитектуры системы выбрана монолитная организация с явно выраженной внутренней модульностью, формализованная в нотации C4~Model. Разработана агентно-независимая топология взаимодействия с удалёнными хостами через протоколы SSH/SFTP, исключающая необходимость предварительной установки дополнительного программного обеспечения на целевые серверы.

Реализованы алгоритм динамического проксирования HTTP-трафика с сохранением метода, заголовков и тела запроса, а также алгоритм ограничения интенсивности запросов Fixed Window Counter на основе атомарных операций Redis. Разработан механизм восстановления состояния системы при перезапуске: процедура restore\_state сверяет метаданные реестра с фактическим состоянием операционной системы и автоматически перезапускает завершившиеся Java-процессы. Реализован модуль экспорта метрик в формате Prometheus с двухпроходной конвейерной схемой сборки данных за фиксированное число обращений к Redis.

Развёртывание и тестирование системы выполнены в среде ОС Astra Linux. По итогам пяти функциональных сценариев подтверждена корректность механизмов персистентности данных (AOF-восстановление Redis) и отказоустойчивости (обнаружение аварийного завершения заглушек планировщиком HealthChecker в пределах 30~секунд). Нагрузочное тестирование показало, что прокси-модуль обеспечивает сквозную задержку не более 3.01~мс при нагрузках до 500~RPS, а алгоритм Rate Limiting удерживает точность ограничения с отклонением 0.13\,\% от теоретического значения.

Практическая значимость работы определяется созданием программного инструмента, позволяющего отказаться от зарубежных решений в задачах управления инфраструктурой имитационных сервисов и обеспечивающего интеграцию со стандартными средствами мониторинга. Направлениями дальнейшего развития системы могут стать поддержка сценариев динамической маршрутизации запросов на основе содержимого тела, реализация ролевой модели разграничения доступа к управляющему интерфейсу, а также расширение поддерживаемых типов артефактов за пределы JVM-платформы.